\cotefile{157-3}
\scanpage{312}
\mspage{262}

carré $W$-cocartésien, qui paraît \unsure{inépuisable} : Deux points de vue, et d'une simplicité parfaite.

Cela amène à revoir la notion en question dans $\mathrm{Cat}$, où elle avait l'air très \unsure{astucieuse}, profonde et un peu obscure, \unsure{un brin vertigineuse|vu du vertiguiste|nom de vertiguiste} et tout ça. C'est le moment de se rendre compte que ``ce n'est que ça''. Donc de prouver la

\uline{Proposition 20} Supposons $\mathcal M=\mathrm{Cat}$, $W\unsure{=}\text{un}$ localiseur fondamental satisfaisant $W\,(7\ \mathrm{bis})$.

Considérons un carré commutatif dans $\mathrm{Cat}$

\diagram{\deleted{\(Q=\left(\begin{tikzcd}[ampersand replacement=\&,column sep=small,row sep=small]
X_{\unsure{1}}\arrow[r,leftrightarrow] \& X_3\arrow[d,leftrightarrow]\\
S\arrow[u,"{\unsure{p_1|p'_1}}"]\arrow[r,leftrightarrow,"{\unsure{p'_2|p_2}}"'] \& X_2
\end{tikzcd}\right).\)}
\qquad \(Q=\left(\begin{tikzcd}[ampersand replacement=\&,column sep=small,row sep=small]
S\arrow[r,dash]\arrow[d] \& X_1\arrow[d]\\
X_2\arrow[r,dash] \& X_3
\end{tikzcd}\right).\)}

Pour qu'il soit $W$-cartésien au sens de la théorie \deleted{\unsure{générale}} développée ici (pour des catégories de modèles générales), il f. et s. que

\diagram{\(\varphi_Q:\left(\overbrace{\int\ \deleted{\substack{\unsure{X_1}\\\unsure{X_2}}}\begin{tikzcd}[ampersand replacement=\&,column sep=tiny,row sep=tiny]
\& X_1\arrow[dl]\\ S\&{}\\ \& X_2\arrow[ul]
\end{tikzcd}}^{\unsure{\mathscr X_Q|X_Q|\mathcal X_Q}}\right)\longrightarrow \deleted{\unsure{D}}X_3\)}

soit dans $W$. (NB L'intérêt de ce critère est qu'il est ``intrinsèque'' : ne faisant pas appel à une factorisation des $p_i$ en des \unsure{fl} en

\diagram{\(\deleted{\unsure{p_i}\longrightarrow\widetilde X_i}\quad\begin{tikzcd}[ampersand replacement=\&,column sep=large]
S\arrow[r,"{\in\mathrm{Cof}\deleted{\unsure{\cap W}}}"'] \& \widetilde X_i\arrow[r,"{\in W}"'] \& X_i
\end{tikzcd}.\unsure{\}}\)}

Cf. applications intéressantes dans section V 3.4. p. 146 ff \unsure{(th.|th.} 4 et corollaires.)
% archive sheet number: 312
\scanpage{313}
\mspage{263}

\uline{Dém} Considérons une factorisation de $S\longrightarrow X_1$ en

\diagram{\(\begin{tikzcd}[ampersand replacement=\&]
S\arrow[r,"i_1","{\unsure{\in|\mathrm C}}"']\&\widetilde X_1\arrow[r,dash,"p_1"]\&X_1
\end{tikzcd},\)}

avec $i_1$ imm ouverte de Dwyer, $p_1\in W$. Donc \unsure{on a} un carré comm.

\diagram{\(\widetilde Q=\left(\begin{tikzcd}[ampersand replacement=\&,column sep=small,row sep=small]
S\arrow[r,hook,"{\mathrm{Dw}}"]\arrow[d]\&\widetilde X_1\arrow[d]\arrow[r,"{\deleted{\unsure{p_1}}}"]\arrow[r,dash,decorate,decoration={zigzag,segment length=2pt,amplitude=1pt}]\&\deleted{\unsure{X_1}}\arrow[d,"{\deleted{\unsure{1}}}"]\arrow[d,dash,decorate,decoration={zigzag,segment length=3pt,amplitude=2pt}]\\
X_2\arrow[r]\&\deleted{\unsure{X_1}}X_3\arrow[r,"{\deleted{\unsure{\leftarrow}}}"]\&X_3
\end{tikzcd}\right)\)}

et un hom $\widetilde Q\xrightarrow{\alpha}\overset{\deleted{\unsure{\sim}}}{Q}$ qui est l'identité sur \inserted{les} trois \deleted{\unsure{des}} sommets distincts de $\widetilde X_1$, et qui est $p_1$ en $\widetilde X_1$.

\unsure{On} en déduit un diagramme comm.

\diagram{\(\begin{tikzcd}[ampersand replacement=\&,column sep=small,row sep=small]
\unsure{\mathscr X|\mathcal X}_Q\arrow[rr,"{\mathscr X_\alpha=\lambda}"]\arrow[dr,"{\varphi_Q}"']\&{}\&\mathscr X_{\widetilde Q}\arrow[dl,dash,"{\varphi_{\widetilde Q}}"]\\
\&X_3\&{}
\end{tikzcd},\)}

et on a $\lambda\in W$ -- (car \deleted{il \unsure{faut}} c'est un hom \unsure{$\Phi$-fibré|$\Phi$-fibr.} de catégories fibrées sur $\Phi$, qui \deleted{\unsure{induit des}} une W-équi. sur chacune des fibres).

% slanted
\margin{On utilise ici $W(3)$ (forme $\Phi$ forte), i.e. $\mathrm{Loc}(3\ \unsure{bis})$ \unsure{ou} -- \unsure{f}, car \uline{cartésiens}, \unsure{de|ou} \uline{\unsure{XVI}}. Mais \unsure{je pense|j'en passe}, cf plus bas.}

Donc $\varphi_Q\in W\Longleftrightarrow\varphi_{\widetilde Q}$ \unsure{tel}, d'autre part il est \deleted{\unsure{pour prouver}}\inserted{clair que} $Q$ est $W$-cocart. ssi $\widetilde Q$ l'est. Donc \unsure{SPS} que $S\hookrightarrow X_1$ est imm. ouverte (de Dwyer). Appliquant ceci
% archive sheet number: 313
\scanpage{314}
\mspage{264}

en changeant les rôles de $X_1$ et $X_2$, on voit qu'on peut supposer de plus que $S\mathbin{\unsure{-|\to}}X_2$ est également imm. \inserted{ouv.} de Dwyer.

Considérons alors

\diagram{\(\begin{tikzcd}[ampersand replacement=\&,column sep=large,row sep=large]
S\arrow[r,arrows={Hooks[right]-}]\arrow[d,hook]\arrow[dr,phantom,"{Q'}" description]\&X_1\arrow[d,hook]\&{}\\
X_2\arrow[r,arrows={Hooks[right]-}]\&X_1\amalg_S X_2\arrow[r,"{\unsure{\pi|\Pi}}"]\&X_3\\
{}\&\unsure{\mathscr X|\mathcal X}_Q\arrow[u,"{\varphi_{Q'}}"]\arrow[ur,"{\varphi_Q}"']\&{}
\end{tikzcd}.\)}

On a $\varphi_Q=\pi\varphi_{Q'}$, d'autre part on sait que $\varphi_{Q'}\in W$ (cor. 3 p. 149). Donc

\[
\pi\in W\Longleftrightarrow\varphi_Q\in W,
\]

or $\pi\in W$ signifie que $Q$ est $W$-cocartésien, qed.
% archive sheet number: 314
\scanpage{315}
\mspage{265}

% slanted
\margin{Réflexions \unsure{provisoires}, \unsure{supplantées} par \uline{XVI}.}

\inserted{(\S 8)} \uline{Appendice} (Retour sur les axiomes $W(7)$, $W(7\ \mathrm{bis})$ (pages 75, 76).

On a vu (prop. 6 p. 78) que $\deleted{\unsure{W(6)}}W(7\ \mathrm{bis})$ implique $W(7)$. Montrons une réciproque partielle :

\uline{Proposition} Supposons $W(7\ \mathrm{bis})$ (axiome \unsure{des|du} \unsure{carré|carrés} \unsure{cocartésien|cocartésiens}) satisfait, et plaçons nous sous les conditions de l'énoncé de $W(7)$, \deleted{\illegible{4}}\inserted{et plus généralement} dans le cas d'un diagramme cubique

\diagram{\(\begin{tikzcd}[ampersand replacement=\&,column sep=small,row sep=small]
{}\&X_2\arrow[rr]\arrow[dd,dashed,"{f_2}" near end]\&{}\&X_3\arrow[dd,"{f_3}"]\\
X_0\arrow[ur,leftrightarrow]\arrow[rr,leftrightarrow]\arrow[dd,"{f_0}"']\&{}\&X_{\unsure{2|1}}\arrow[ur]\arrow[dd,"{f_{\unsure{1|2}}}"']\&{}\\
{}\&Y_2\arrow[rr,dashed]\&{}\&Y_3\\
Y_0\arrow[ur]\arrow[rr]\&{}\&Y_{\unsure{1|2}}\arrow[ur]\&{}
\end{tikzcd},\)}

où les deux carrés $X$ et $Y$ (faces sup. et inf. du cube) sont $W$-cocartésiens. Considérons les conditions suivantes sur ce diagramme

\[
\begin{aligned}
&\text{a)}\ f_0\in W\\
&\text{b)}\ f_1\text{ et }f_2\in W\\
&\text{c)}\ f_3\in W.
\end{aligned}
\]

Alors on a les implications

\[
a\mathbin{\&}b\Longrightarrow c,\qquad b\mathbin{\&}c\Longrightarrow a.
\]

De plus si $W$ est stable par facteurs directs (p. ex. $W$ fortement saturé), alors on a aussi

\[
a\mathbin{\&}c\Longrightarrow b
\]

donc \inserted{la conjonction \unsure{de}} deux \inserted{quelconques} parmi les conditions a) b) c) implique la troisième.

Démonstration. Utilisant la factorisation
% archive sheet number: 315
\scanpage{316}
\mspage{266}

fonctorielle d'une flèche quelconque (ici les $f_i$) en $pi$, avec $i$ imm ouverte de Dwyer, $p\in\unsure{W_1|W,}$ (p. 229) on est ramené au cas où les $f_i$ sont des imm. ouvertes de Dwyer. \unsure{So} Mais pour \deleted{\unsure{toute|de t}} flèche

% slanted
\margin{C'est un peu vif, car il n'est pas clair que $X'_1\amalg_{X_0}X'_2\longrightarrow X_1\amalg_{X_0}X_2$ soit $\in W$, sans utiliser déjà \deleted{\unsure{la propriété}} \inserted{\unsure{disons}} \unsure{l'un} des trois \unsure{ingrédients} \unsure{$\bigcirc$} de $W(7)$ (savoir $f_0,f_1,f_2\in W\Longrightarrow f_3\in W$). Il semble donc que notre argument ici, \unsure{n'aurait lieu} que moyennant la mise \unsure{hors} \deleted{\unsure{déjà}} \unsure{faite} d'une partie des \unsure{cas admis} \unsure{\ldots} Mais il y a le ``\uline{lemme de Brown} '', cf \uline{XVI} \unsure{\S 3|\S 1}, \unsure{cet argument|ce qui suit} p. \unsure{2|1} \unsure{ff}.}

\inserted{$g:S\mathbin{\unsure{-|\to}}T$} qui est dans $\mathrm{Cof}_W$, soit $Z(g)$ \unsure{son|sa} \unsure{mapping-cone}

\diagram{\(\begin{tikzcd}[ampersand replacement=\&,column sep=small,row sep=small]
S\arrow[rr,"g"]\arrow[dr]\&{}\&T\arrow[dr,"{\alpha(g)}"]\&{}\\
{}\&e\arrow[rr,"{\overline g}"]\&{}\&Z(g)
\end{tikzcd},\)}

\[
\alpha(g):T=\mathrm{but}(g)\longrightarrow Z(g)
\]

la flèche canonique. Celle-ci est fonctorielle en $g$, de sorte que le diagramme cubique donné se prolonge en un double cube

\diagram{\(\begin{tikzcd}[ampersand replacement=\&,row sep=small,column sep=small]
{}\&X_2\arrow[rr]\arrow[dd,dashed,"{f_2}" near end]\&{}\&X_3\arrow[dd,"{f_3}"]\\
X_0\arrow[ur]\arrow[rr]\arrow[dd,"{f_0}"']\&{}\&X_{\unsure{2|1}}\arrow[ur]\arrow[dd,"{f_1}"']\&{}\\
{}\&Y_2\arrow[rr,dashed]\arrow[dd,dashed,"{\alpha_2}" near start]\&{}\&Y_3\arrow[dd,"{\alpha_3}"]\\
Y_0\arrow[ur]\arrow[rr]\arrow[dd,"{\alpha_0}"']\&{}\&Y_1\arrow[ur]\arrow[dd,"{\alpha_1}"']\&{}\\
{}\&Z_2\arrow[rr,dashed]\&{}\&Z_3\\
Z_0\arrow[ur]\arrow[rr,line width=0.9pt]\&{}\&Z_1\arrow[ur]\&{}
\end{tikzcd}\)}

avec \deleted{$\unsure{Z}$} $Z_i=Z(f_i)$, $\alpha_{\unsure{i}}=\alpha(f_i)$ pour $i\in\{0,1,2,3\}$.

Notons que $W(7\ \mathrm{bis})$ implique que si $g\in\mathrm{Cof}_W$,
% archive sheet number: 316
\scanpage{317}
\mspage{267}

on a

\[
\deleted{\unsure{\alpha}}g\in W\Longleftrightarrow\overline{\deleted{\unsure{Z}}g}\in W\quad\text{i.e.}\quad Z(g)\ W\text{-asph.},
\]

en particulier

\[
\deleted{\unsure{\alpha}}f_i\in W\Longleftrightarrow Z_i\ \deleted{\unsure{d}}W\text{-asph.}
\]

Ainsi, l'énoncé à prouver sur le cube initial se ramène à \deleted{\unsure{celle}} un énoncé \deleted{\uline{Corollaire} 1. Moyennant $W(7\ \mathrm{bis})$, si on a un diagramme} sur le carré

\diagram{\(\begin{tikzcd}[ampersand replacement=\&,row sep=small,column sep=small]
{}\&Z_1\arrow[dr]\&{}\\
Z_0\arrow[ur]\arrow[dr,dash]\&{}\&Z_3\\
{}\&Z_2\arrow[ur]\&{}
\end{tikzcd},\)}

savoir que les trois conditions

\[
\begin{aligned}
&\text{a}')\quad Z_0\ \deleted{\unsure{d}}W\text{-asph.}\\
&\text{b}')\quad Z_1\text{ et }Z_2\quad W\text{-asph.}\\
&\text{c}')\quad Z_3\ W\text{-asph.}
\end{aligned}
\]

donnent lieu aux implications

\[
a'\mathbin{\&}b'\Longrightarrow c',\qquad b'\mathbin{\&}c'\Longrightarrow a',
\]

et que si $W$ est stable par facteurs directs, alors

\[
a'\mathbin{\&}c'\Longrightarrow b'.
\]

C'est à peu de choses près le cor. de la page 76, complété à la page 77 (corollaire), à cela près qu'on y suppose les flèches des imm. ouvertes, et le carré cartésien \deleted{et \unsure{cocartésien}}\inserted{(\unsure{donc|et} aussi cocartésien)} dans $\mathrm{Cat}$, \uline{i.e.} $Z_3=Z_1\cup Z_2$, $Z_0\mathrel{\unsure{=|>}}Z_1\cap Z_2$. Mais
% archive sheet number: 317
\scanpage{318}
\mspage{268}

ce qui importe pour la validité de cet énoncé, c'est (on s'en doute) que $Z_{\unsure{*|\bullet}}$ soit $W$-cocartésien. Donc il faut d'abord prouver ceci

\uline{Lemme} Si le \deleted{\unsure{cub}}diagramme cubique p. 265 est tel que les faces $X,Y$ soient $W$-cocartésiens, et \deleted{\unsure{$f_i$}} les $f_i$ ($i\in\{0,1,2,3\}$) des $W$-cofibrations, alors le carré $Z$ \deleted{qui s'en \unsure{déduit} en prenant} des cofibres des $f_i$ est lui aussi $W$-cocartésien.

Ce lemme résulte de la théorie du ``déploiement des diagrammes cubiques'', \unsure{relatif} \deleted{\unsure{une}} à une catégorie de modèles $(\mathcal M,W)$ admettant une théorie des cofibres $W$-homotopiques, qu'il me faudra bien écrire ! (Mais l'existence \inserted{\deleted{\unsure{pour}} $(\mathrm{Cat},W)$} d'une telle théorie des cofibres $W$-homotopiques, et par suite de celle du déploiement des diagrammes dans $\mathcal M$ de type $\unsure{\Delta_1^{I}|\Delta^{1\,I}}$, $I$ ens. d'indices \unsure{finis|fini}, ne dépend que de $W(7\ \mathrm{bis})$ -- on n'a pas \unsure{à} faire appel à $W(7)$.)
% archive sheet number: 318
\scanpage{319}
\mspage{269}

Ce ``déploiement'' est un diagramme \inserted{infini} dont toutes les suites infinies, dans les trois directions des axes \unsure{de|des} coordonnées, sont des ``suites de suspension''. Je représente la portion utile du déploiement

\diagram{\(\begin{tikzcd}[ampersand replacement=\&,row sep=tiny,column sep=tiny]
{}\&{}\&{}\&{}\&{}\&{}\&\unsure{U|u}\arrow[dd,"{\mu}"]\\
{}\&X_{\deleted{\unsure{3}}2}\arrow[ur]\arrow[rr]\arrow[dd,dashed]\&{}\&X_3\arrow[ur]\arrow[rr,dash]\arrow[dd,dashed]\&{}\&T\arrow[ur]\arrow[dd]\&{}\\
X_0\arrow[ur]\arrow[rr]\arrow[dd]\&{}\&X_1\arrow[ur,dash]\arrow[rr]\arrow[dd]\&{}\&S\arrow[ur,"{\lambda}"]\arrow[dd]\&{}\&U'\arrow[dd]\\
{}\&Y_2\arrow[ur,shorten >=5pt]\arrow[rr,dashed]\arrow[dd,dashed]\&{}\&Y_3\arrow[ur,dashed,shorten >=6pt]\arrow[rr,dashed]\arrow[dd,dashed]\&{}\&\deleted{\unsure{\overline T}}T'\arrow[ur]\arrow[dd]\&{}\\
Y_0\arrow[ur,dash]\arrow[rr]\arrow[dd]\&{}\&Y_1\arrow[ur]\arrow[rr]\arrow[dd]\&{}\&S^{\unsure{\prime}}\arrow[ur,"{\lambda'}"]\arrow[dd]\&{}\&U''\\
{}\&Z_2\arrow[ur,shorten >=5pt]\arrow[rr,dashed]\&{}\&Z_3\arrow[dl]\arrow[ur,dashed,shorten >=6pt]\arrow[rr,dashed]\&{}\&T''\arrow[ur,line width=1pt]\&{}\\
Z_0\arrow[ur]\arrow[rr]\&{}\&Z_1\arrow[rr,dash]\&{}\&S''\arrow[ur,dash,"{\lambda''}"]\&{}\&{}
\end{tikzcd}\)}

La condition que $X$ soit $W$-cocartésien s'exprime par $\lambda\in W$, ou encore par $U\deleted{\unsure{\in}}\in\mathrm{Asph}_W$ (car $U$ s'identifie à \deleted{\unsure{$Z$|la}} \unsure{cofibre} homotopique $Z(\lambda)$ de $\lambda$ \unsure{\ldots}), de $=$ $Y$ $W$-cocartésien s'exprime par $U'\in\mathrm{Asph}_W$, et $Z$ $W$-cocartésien par $U''\in\mathrm{Asph}_W$. Or $U''\sim Z(\mu)$, où $\mu:U\longrightarrow U'$ est dans $W$ si $U,U'$ sont $W$-asphériques, donc $U''=Z(\mu)$ est alors $W$-asph. Cela prouve le lemme.

Cela montre que, moyennant $W(7\ \mathrm{bis})$ (en fait, on a utilisé seulement la partie ``universelle'' de $W(7\ \mathrm{bis})$, savoir que dans

\diagram{\(\deleted{\unsure{\#}}\quad\begin{tikzcd}[ampersand replacement=\&,row sep=tiny,column sep=small]
T_0\arrow[r]\arrow[d,"i"']\&T_2\arrow[d,"{i'}"]\\
T_{\unsure{1}}\arrow[r]\&T_{\unsure{3}}
\end{tikzcd}\)}

diagr. cartésien d'imm. ouvertes, $i\in W\Longrightarrow i'\in W$ (au lieu de $\Longleftrightarrow$),
% archive sheet number: 319
\scanpage{320}
\mspage{270}

la validité de $W(7)$ se ramène au cas où le carré $Y$ est le carré trivial

\diagram{\(\begin{tikzcd}[ampersand replacement=\&,column sep=tiny,row sep=tiny]
{}\&e\arrow[dr]\&{}\\e\arrow[ur]\arrow[dr]\&{}\&e\\{}\&e\arrow[ur]\&{}
\end{tikzcd},\)}

$e$ l'objet final de $\mathrm{Cat}$.

On peut maintenant conclure : \deleted{\unsure{Si}} On a

\diagram{\(\begin{tikzcd}[ampersand replacement=\&,column sep=small,row sep=small]
{}\&Z_1\arrow[dr,"{j_2}"]\&{}\\Z_0\arrow[ur,dash,"{i_1}"]\arrow[dr,"{i_2}"']\&{}\&Z_3\\{}\&Z_2\arrow[ur,"{j_1}"']\&{}
\end{tikzcd}\)}

(moyennant la partie directe de $W(7\ \mathrm{bis})$)

\[
a'\mathbin{\&}b'\Longrightarrow c'
\]

car si $Z_0,Z_1\deleted{\unsure{,Z_2}}$ sont $\in\mathrm{Asph}_W$, alors $i_1\in W$, donc $j_1\in W$, et si de plus $Z_2\in\mathrm{Asph}_W$, on a donc $Z_3\in\mathrm{Asph}_W$.

Si on a aussi la partie inverse \inserted{(``\unsure{inhabituelle|inconditionnelle}'')} de $W(7\ \mathrm{bis})$, alors

\[
b'\mathbin{\&}c'\Longrightarrow a',
\]

car si $Z_1$ et $Z_3$ sont \unsure{$W$-asph|$U$-asph}, alors $j_2\in W$, donc $i_2\in W$, et si de plus $Z_2$ est $W$ asph, on conclut $Z_0$ $W$-asph.

Enfin, supposons que $W$ soit \uline{stable par facteurs directs}, et $W(7\ \mathrm{bis})$ satisfait sous forme directe. Je dis que alors

\[
c'\mathbin{\&}a'\Longrightarrow b'
\]

\[
\text{i.e.}\qquad Z_0\text{ et } Z_3\in\mathrm{Asph}_W\Longrightarrow Z_1\text{ et }Z_2\in\mathrm{Asph}_W.
\]

Factorisant $i_1$ et $i_2$ en

\diagram{\(\begin{tikzcd}[ampersand replacement=\&,column sep=small,row sep=tiny]
Z_0\arrow[r,dash,"{i'_1}"]\&Z'_1\arrow[r,"{\alpha_1}"]\&Z_1,\\
Z_0\arrow[r,dash,"{i'_2}"]\&Z'_2\arrow[r,dash,"{\alpha_2}"]\&Z_2,
\end{tikzcd},\)}

avec $i'_1,i'_2$ des imm. ouvertes de Dwyer, et $\alpha_1,\alpha_2\in W$, on est ramené au cas où
% archive sheet number: 320
\scanpage{321}
\mspage{\deleted{\unsure{270}}271}

$i_1,i_2$ sont des \uline{imm. ouvertes} (de Dwyer)\inserted{\deleted{\unsure{asphériques|symétriques}}} et où le carré est cocartésien (ce qui implique qu'il est $W$-cocartésien). Faisant alors le \unsure{cochgt.} de base par $Z_0\xrightarrow{p_0}Z'_0=e$, on trouve un \deleted{carré} hom de carrés cocartésiens

\diagram{\(\begin{tikzcd}[ampersand replacement=\&,column sep=small,row sep=small]
{}\&Z_1\arrow[dr]\arrow[ddd,dashed,"{p_1}" near start]\&{}\\
Z_0\arrow[ur]\arrow[dr]\arrow[ddd,"{\deleted{\unsure{p}}p_0}"']\&{}\&Z_3\arrow[ddd,"{p_3}"]\\
{}\&Z_2\arrow[ur]\arrow[ddd,dash,"{p_2}" near end]\&{}\\
{}\&Z'_1\arrow[dr,dashed]\&{}\\
e=Z'_0\arrow[ur]\arrow[dr]\&{}\&Z'_3\\
{}\&Z'_2\arrow[ur]\&{}
\end{tikzcd}\)}

où toutes les faces du cube sont $W$-cocartésiennes, et comme $p_0\in W\unsure{,}$ ($Z_0$ \unsure{étant} $W$-asph) les autres flèches verticales $p_1,p_2,p_3$ sont $\in W$, donc le carré $Z'_{\unsure{*|\bullet}}$ satisfait aussi $Z'_0$ et $Z'_3$ asph (et de plus $Z'_0=e$), et il faut prouver $Z'_1,Z'_2$ $W$-asph (d'où idem pour $Z_1,Z_2$). Donc, on est ramené au cas particulier où $Z_0=e$.

Mais dans ce cas il existe une rétraction \inserted{(triviale !)} de $Z_1$ sur $Z_0$, donc aussi de $Z_3$ sur $Z_2$. Ainsi, $Z_2$ est un facteur
% archive sheet number: 321
\scanpage{322}
\mspage{272}

direct d'un objet $W$-asphérique, ce qui implique, par l'hypothèse sur $W$, que $Z_2$ est $W$-asphérique. \inserted{Par} symétrie, $Z_1$ l'est aussi. On a gagné !

\uline{NB} En fait, ces arguments sont généraux, et donnent ceci :

\uline{Théorème} Soit $(\mathcal M,W)$ une cat. de modèles, adme\deleted{\unsure{tto}}ttant une théorie de la cofibre $W$-homotopique\({}^{*}\) (i.e. on a une factorisation de \deleted{\unsure{d}}tte flèche $f:X\longrightarrow Y$ en $pi$, avec $i\in\mathrm{Cof}_W$, $p\in W$.) \deleted{\unsure{Et}} On a donc une notion de carré $W$-\unsure{cocartésiens|cocartésien} dans $\mathcal M$. Soit alors $f:X\mathbin{\unsure{-|\to}}Y$

% margin note written sideways
\margin{$*$ Mais on ne sait \unsure{vérifier} l'existence d'une telle théorie dans $(\mathrm{Cat},W)$ que si on admet la forme faible de $W(7)$, i.e. $f_0,f_1,f_2\in W\Longrightarrow f_3\in W$, ou quelque \unsure{cas particulier|permutation} de celle-ci. Voir le \uline{lemme de Brown} dans la suite. Cf \uline{XV} \S 3.}

% slanted
\margin{à vérifier que la factorisation suffit, $=$ si tous les objets ne sont pas fibrants sur $e$. Mais on peut alors remplacer $\mathcal M$ par $\mathcal M_{\unsure{f}}$, et développer le formalisme des fibres homotopiques en conséquence\ldots}

\diagram{\(\begin{tikzcd}[ampersand replacement=\&,row sep=small,column sep=small]
X_0\arrow[rr]\arrow[dr,dash]\arrow[dd,"{f_0}"']\&{}\&X_2\arrow[dr]\arrow[dd,dashed,"{f_2}" near end]\&{}\\
{}\&X_1\arrow[rr]\arrow[dd,"{f_1}"']\&{}\&X_3\arrow[dd,"{f_3}"]\\
Y_0\arrow[rr,dashed]\arrow[dr,dash]\&{}\&Y_2\arrow[dr]\&{}\\
{}\&Y_1\arrow[rr]\&{}\&Y_3
\end{tikzcd}\)}

un hom. d'un carré $W$-cocartésien \inserted{$X$} dans un autre $Y$ $f=(f_0,f_1,f_2,f_3)$.

Considérons les conditions sur $f$

\[
\begin{aligned}
&\text{a)}\ f_0\mathrel{\unsure{\in|+}}W\\
&\text{b)}\ f_1\text{ et }f_2\in W\\
&\text{c)}\ f_3\in W.
\end{aligned}
\]

Alors :

\[
1^\circ)\quad\deleted{\unsure{f}}a\deleted{\unsure{\&}}\mathbin{\&}b\Longrightarrow c.
\]
% archive sheet number: 322
\scanpage{323}
\mspage{273}

$\deleted{\unsure{b}}2^\circ)$ \deleted{Si dans $(\mathcal M,W)$, \unsure{on a}} \inserted{Supposons que $\mathcal M$ ait un objet final $e$. Alors} l'implication

\[
b\mathbin{\&}c\Longrightarrow a
\]

pour \uline{tout} hom $f:X\longrightarrow Y$ de carrés $W$-cocartésiens équivaut à l'axiome suivant sur $W$ :

\unsure{(}\uline{Axiome des carrés \inserted{$W$-}cocartésiens conservatifs} : Si

\diagram{\(\begin{tikzcd}[ampersand replacement=\&,column sep=small,row sep=small]
{}\&Z_1\arrow[dr,dash,"{j_2}"]\&{}\\Z_0\arrow[ur,"{i_1}"]\arrow[dr,"{i_2}"']\&{}\&Z_3\\{}\&Z_2\arrow[ur,"{j_1}"']\&{}
\end{tikzcd}\)}

% slanted
\margin{Cf argument page 78 qui montre que l'implication $b\mathbin{\&}c\Longrightarrow a$ entraîne l'axiome conservatif}

$W$-cocartésien, alors $j_1\in W\Longrightarrow i_1\in W$ (donc en fait $i_1\in W\Longleftrightarrow j_1\in W$ et $i_2\in W\Longleftrightarrow j_2\in W$).

Cet axiome est équivalent à la conjonction des deux axiomes \unsure{suivant|suivants}

\uline{Cons. 1} Si on a un carré \deleted{\unsure{$W$-}}cocartésien

\diagram{\(\begin{tikzcd}[ampersand replacement=\&,column sep=small,row sep=small]
\deleted{\unsure{X}}S\arrow[r,"f"]\arrow[dr,dash]\&\deleted{\unsure{Z}}T\arrow[dr,"{\deleted{\unsure{\mathrm{ou}}}}"]\&{}\\{}\&e\arrow[r,dash]\&U
\end{tikzcd}\qquad\unsure{(-}\)}

\deleted{(\unsure{où on peut supposer}} \inserted{avec} $f\in\mathrm{Cof}_W$, \deleted{et le carré cartésien alors} alors $U\in\mathrm{Asph}_W\Longrightarrow f\in W$ (donc $U\in\mathrm{Asph}_W\Longleftrightarrow f\in W$)

\uline{Cons. 2} Si on a un carré cartésien,

\diagram{\deleted{avec}\quad\(\deleted{\unsure{S}\longrightarrow}\begin{tikzcd}[ampersand replacement=\&,row sep=small,column sep=small]
{}\&Z_1\arrow[dr,dash]\&{}\\Z_0\arrow[ur,dash,"{i_1}"]\arrow[dr,dash,"{i_2}"']\&{}\&Z_3\\{}\&Z_2\arrow[ur,dash]\&{}
\end{tikzcd}\)}

avec $i_1,i_2\in\mathrm{Cof}_W$, alors si $Z_1,Z_2,Z_3$ sont $W$ asph, $Z_0$ l'est aussi.
% archive sheet number: 323
\scanpage{324}
\mspage{274}

$3^\circ)$ Supposons que tout\deleted{\unsure{s}} facteur direct d'un objet $W$-asphérique soit $W$-asphérique, ce qui est le cas notamment si $W$ est \deleted{\unsure{lui}} $=$ stable par facteurs directs, a fortiori si $W$ est fortement saturé. \inserted{Supposons de plus que $\mathcal M$ ait \unsure{objet|l'objet} final.} Alors on a (sous les conditions générales du début)

\[
c\mathbin{\&}a\Longrightarrow b
\]

i.e. si $f_0,f_3$ sont $\in W$, de $=$ $f_1,f_{\unsure{2}}$.

Cette dernière \deleted{\unsure{assertion}}\inserted{\unsure{propriété}} peut se déployer ainsi

\uline{Corollaire 1} \unsure{Sans|Sous} hyp. suppl. $3^\circ)$ sur $W$, la validité de $c\mathbin{\&}a\Longrightarrow b$ (pour toute flèche $f$ entre carrés $W$-cocartésiens) équivaut à \unsure{chacun|chacune} des deux \inserted{autres} axiomes suivants sur $W$ (qui sont donc équivalents) :

1) Si on a un carré \deleted{cartésien} $W$-cocartésien

\diagram{\(\begin{tikzcd}[ampersand replacement=\&,row sep=small,column sep=small]
{}\&Z_1\arrow[dr]\&{}\\Z_0\arrow[ur]\arrow[dr]\&{}\&Z_3\\{}\&Z_2\arrow[ur]\&{}
\end{tikzcd},\)}

alors $Z_0$ et $Z_3$ $W$-asph $\Longrightarrow\unsure{Z_1|Z}$ et $\unsure{Z_2|Z}$ $W$-asphériques

1$'$) Comme 1), avec $Z_0=e$ (dans le cas où on suppose que $\mathcal M$ admet objet final $e$).
% archive sheet number: 324
\scanpage{325}
\mspage{275}

\uline{Cor. 2} Supposons que $W$ satisfasse l'axiome des carrés $W$-cocartésiens conservatifs, que tout facteur direct d'un objet $W$-asphérique soit $W$-asphérique (p. ex. $W$ stable par facteurs directs), enfin que $\mathcal M$ admette un objet final. Alors pour \unsure{tt} \deleted{\unsure{$f$}} hom $f:X\mathbin{\unsure{-|\to}}Y$ \unsure{des|de} carrés $W$-cocartésiens, \inserted{la conjonction de} deux parmi les conditions a) b) c) de p. 272 entraîne la troisième.
% archive sheet number: 325
\scanpage{326}

\handwritten{I) Contrexemples relatifs à \(\underline{\overline{\mathrm{XII}}}\) (\unsure{ex|ex.} d'immersions ouvertes $W_\infty$-asphériques qui ne sont pas des $W_\infty$-cofibrations \hfill $1$--$16$}

\handwritten{II) Recollement d'un ouvert et d'un fermé comme \unsure{$W_\infty$-sommes|$W_\infty$-somme} amalgamée\unsure{s} (?)}

% printed form
% Reverse-side print visible through the sheet, read in the direction of the type.
\begin{flushleft}\ttfamily
\unsure{JOB|J0B} \unsure{1054}\\
00000020\\
00000030\unsure{.}\\
\unsure{IT=SYSDA,|T=SYSDA,}\qquad 00000040\\
\unsure{(3,1),RLSE)|(3,1);RLSE)}\qquad 00000050\\
00000060\\
00000070\\
00000080\\
00000090\\
00000100\\
00000110\\
00000140
\end{flushleft}
% archive sheet number: 326
